\section{Irreducibility of the dif\/ference Galois group}
\label{sec4}
We let
\begin{gather}
\label{equation d ordre 2}
\phi^{2}(y) + a \phi(y)+by=0 \qquad \text{with} \quad a \in \operatorname{M}_{\Lambda} \quad \text{and} \quad b\in
\operatorname{M}_{\Lambda}^{\times}
\end{gather}
be a~dif\/ference equation of order $2$ with coef\/f\/icients in $\operatorname{M}_{\Lambda}$ and we denote~by
\begin{gather*}
\phi Y=AY \qquad \hbox{with} \quad A=
\begin{pmatrix}
0&1
\\
-b&-a
\end{pmatrix}
\in \operatorname{GL}_{2}(\operatorname{M}_{\Lambda})
\end{gather*}
the associated dif\/ference system.
For the notations $\operatorname{M}_{\Lambda}$,~$\phi$,~$K$, etc, we refer to Sections~\ref{sec2} and~\ref{sec3}.

We let $G \subset \operatorname{GL}_{2}(\mathbb C)$ be the dif\/ference Galois group over $(K,\phi)$ of
equation~\eqref{equation d ordre 2}.
According to Corollary~\ref{coro1},~$G$ is an algebraic subgroup of $\operatorname{GL}_{2}(\mathbb C)$ such that the
quotient $G/G^{\circ}$ of~$G$ by its identity component $G^{\circ}$ is cyclic.
A~direct inspection of the classif\/ication, up to conjugation, of the algebraic subgroups of
$\operatorname{GL}_{2}(\mathbb C)$ given in~\cite[Theorem~4]{NvdPT} shows that~$G$ satisf\/ies one of the following
properties:
\begin{itemize}
\itemsep=0pt
\item The group~$G$ is reducible (i.e., conjugate to some subgroup of the group of upper-triangular matrices in
$\operatorname{GL}_{2}(\mathbb C)$).
If~$G$ is reducible, we distinguish the following sub-cases:
\begin{itemize}\itemsep=0pt
\item the group~$G$ is completely reducible (i.e., is conjugate to some subgroup of the group of diagonal matrices
in $\operatorname{GL}_{2}(\mathbb C)$);
\item the group~$G$ is not completely reducible.
\end{itemize}
\item The group~$G$ is irreducible (i.e., not reducible) and imprimitive (see Section~\ref{sec5} for the def\/inition).
\item The group~$G$ is irreducible and is not imprimitive, and, in this case, there exists an algebraic subgroup~$\mu$
of $\mathbb C^{\times}$ such that $G= \mu \operatorname{SL}_{2}(\mathbb C)$.
Therefore, $G= \{M \in \operatorname{GL}_{2}(\mathbb C) \, \vert \, \det(M) \in H\}$ where $H=\det(G) \subset \mathbb
C^{\times}$.
In order to determine~$H$, one can use the fact that $H=\det(G)$ is the dif\/ference Galois group of $\phi y = (\det A) y
= by$ (this follows for instance from Tannakian duality~\cite[Section~1.4]{VdPS97}).
\end{itemize}

Our f\/irst task, undertaken in the present section, is to study the reducibility of~$G$.
The imprimitivity of~$G$ will be considered in~Section~\ref{sec5}.

\subsection{Riccati equation and irreducibility}

The non linear dif\/ference equation
\begin{gather}
\label{eq4}
(\phi(u)+a) u=-b
\end{gather}
is called the Riccati equation associated to equation~\eqref{equation d ordre 2}.
A~straightforward calculation shows that~$u$ is a~solution of this equation if and only if $\phi-u$ is a~right factor of
$\phi^{2}+a\phi+b$, whence its link with irreducibility.

In what follows, we denote by $I_{2}$ the identity matrix of $\operatorname{GL}_2(\mathbb C)$.

\begin{Lemma}
\label{lem1}
The following statements hold:
\begin{enumerate}\itemsep=0pt
\item[$1.$] If~\eqref{eq4} has one and only one solution in~$K$ then~$G$ is reducible but not completely reducible.
\item[$2.$] If~\eqref{eq4} has exactly two solutions in~$K$ then~$G$ is completely reducible but not an algebraic subgroup of
$\mathbb C^{\times} I_{2}$.
\item[$3.$] If~\eqref{eq4} has at least three solutions in~$K$ then it has infinitely many solutions in~$K$ and~$G$ is an
algebraic subgroup of $\mathbb C^{\times} I_{2}$.
\item[$4.$] If none of the previous cases occur then~$G$ is irreducible.
\end{enumerate}
\end{Lemma}

\begin{proof}
The proof of this lemma is identical to that of~\cite[Theorem~4.2]{Hen98}, to whom we refer for more details.

(1) We assume that~\eqref{eq4} has one and only one solution $u\in K$.
A~straightforward calculation shows that
\begin{gather*}
\phi(T)AT^{-1}=
\begin{pmatrix}
u&*
\\
0&b/u
\end{pmatrix}
\qquad \text{for} \quad T:=
\begin{pmatrix}
1-u&1
\\
-u&1
\end{pmatrix}
\in\operatorname{GL}_{2}(K).
\end{gather*}
We deduce from this and from Corollary~\ref{coro1} that~$G$ is reducible.

Moreover, if~$G$ was completely reducible then, in virtue of Corollary~\ref{coro1}, $\phi(T)AT^{-1}$ would be diagonal
for some ${T:= (t_{i,j})_{1 \leq i,j\leq 2}\in \operatorname{GL}_{2}(K)}$.
Equating the entries of the antidiagonal of $\phi(T)AT^{-1}$ with~$0$, we f\/ind that
$-\frac{t_{21}}{t_{22}},-\frac{t_{11}}{t_{12}}\in K $ are solutions of the Riccati equation.
Since $\det(T) \neq 0$, these solutions are distinct, whence a~contradiction.

(2) Assume that~\eqref{eq4} has exactly two solutions $u_{1},u_{2}\in K$.
We have
\begin{gather*}
\phi(T)AT^{-1}=
\begin{pmatrix}
u_{1}&0
\\
0&u_{2}
\end{pmatrix}
\qquad \text{for}\quad T:=\frac{1}{u_{1}-u_{2}}
\begin{pmatrix}
-u_{2}&1
\\
-u_{1}&1
\end{pmatrix}
\in \operatorname{GL}_{2}(K).
\end{gather*}
We deduce from this and from Corollary~\ref{coro1} that~$G$ is completely reducible.

Moreover, if~$G$ was an algebraic subgroup of $\mathbb C^{\times} I_{2}$ then, according to Corollary~\ref{coro1}, there
would exist $u\in K$ and $T= (t_{i,j})_{1 \leq i,j\leq 2} \in \operatorname{GL}_{2}(K)$ such that
\begin{gather*}
\phi(T)AT^{-1}=u I_{2}.
\end{gather*}
This equality implies that $t_{21}$ and $t_{22}$ are non zero and that, for all $c,d\in \mathbb C$ with
$ct_{2,2}+dt_{1,2} \neq 0$,
\begin{gather*}
-\frac{ct_{21}+dt_{11}}{ct_{22}+dt_{12}}\in K
\end{gather*}
is a~solution of~\eqref{eq4}.
It is easily seen that we get in this way inf\/initely many solutions of the Riccati equation, this is a~contradiction.

(3) Assume that~\eqref{eq4} has at least three solutions $u_{1},u_{2},u_{3}\in K$.
The proof of assertion (2) of the present lemma shows that $\phi Y=AY$ is isomorphic over~$K$ to $\phi Y=
\begin{pmatrix}
u_{i}&0
\\
0&u_{j}
\end{pmatrix}
Y$ for all $1\leq i<j \leq 3$.
Therefore, there exists $T \in \operatorname{GL}_{2}(K)$ such that
\begin{gather*}
\phi(T)
\begin{pmatrix}
u_{1}&0
\\
0&u_{2}
\end{pmatrix}
=
\begin{pmatrix}
u_{1}&0
\\
0&u_{3}
\end{pmatrix}
T.
\end{gather*}
Equating the second columns in this equality, we see that there exists ${f\in K^{\times}}$ such that either
$u_{1}=\frac{\phi f}{f}u_{2}$ or $u_{3}=\frac{\phi f}{f}u_{2}$; up to renumbering, one can assume that the former case
holds true.
It follows that $\phi Y=AY$ is isomorphic over~$K$ to
\begin{gather*}
\phi Y=(u_{1} I_{2}) Y
\end{gather*}
and, according to Corollary~\ref{coro1},~$G$ is an algebraic subgroup of $\mathbb C^{\times} I_{2}$.
We have shown during the proof of statement (2) that this implies that the Riccati equation~\eqref{eq4} has inf\/initely
many solutions in~$K$.

(4) Assume that~$G$ is reducible.
According to Corollary~\ref{coro1}, there exists $T=(t_{i,j})_{1 \leq i,j\leq 2}\in \operatorname{GL}_{2}(K)$ such that
$\phi(T)AT^{-1}$ is upper triangular.
Then $t_{22}\neq 0$ and $-\frac{t_{21}}{t_{22}}\in K$ is a~solution of the Riccati equation~\eqref{eq4}.
This proves claim (4).
\end{proof}

In the proof of the previous lemma, we have shown the following result, which we state independently for ease of
reference.

\begin{Lemma}
\label{lem ric irred bis}
The following properties are equivalent:
\begin{itemize}\itemsep=0pt
\item The Riccati equation~\eqref{eq4} has at least three solutions in~$K$.
\item The Riccati equation~\eqref{eq4} has infinitely many solutions in~$K$.
\item The difference Galois group~$G$ is a~subgroup of $\mathbb C^{\times}I_{2}$.
\item There exist $u \in K^{\times}$ and $T \in \operatorname{GL}_2(K)$ such that $\phi(T)AT^{-1} = u I_{2}$.
\end{itemize}
\end{Lemma}

We shall now state and prove one more lemma.

\begin{Lemma}
\label{lem descente}
Let $\Lambda'' \subset \Lambda'$ be sublattices of~$\Lambda$ such that the quotient $\Lambda'/\Lambda''$ is cyclic.
Assume that there exist $u \in \operatorname{M}_{\Lambda''}^{\times}$ and $T \in
\operatorname{GL}_{2}(\operatorname{M}_{\Lambda''})$ such that
\begin{gather}
\label{et une equation une pour lem}
\phi(T)AT^{-1}=u I_{2}.
\end{gather}
Then, the Riccati equation~\eqref{eq4} has at least two distinct solutions in $\operatorname{M}_{\Lambda'}$.
\end{Lemma}

\begin{proof}
The Galois extension $ \operatorname{M}_{\Lambda''} \vert \operatorname{M}_{\Lambda'} $ is cyclic of order
$k:=[\operatorname{M}_{\Lambda''}: \operatorname{M}_{\Lambda'}]$.
Its Galois group $\operatorname{Gal}(\operatorname{M}_{\Lambda''}|\operatorname{M}_{\Lambda'})$ is generated by the
f\/ield automorphism $\sigma_{1}$ given by $\sigma_{1}(f(z))=f(z+\lambda')$, where $\lambda' \in \Lambda'$ is
a~representative of a~generator of $\Lambda'/\Lambda''$.
Note that the action of $\operatorname{Gal}(\operatorname{M}_{\Lambda''}|\operatorname{M}_{\Lambda'})$ on~$\operatorname{M}_{\Lambda''}$ commutes with the action of~$\phi$.
Applying~$\sigma_{1}$ to equation~\eqref{et une equation une pour lem}, we get
\begin{gather*}
\phi(\sigma_{1} (T))A\sigma_{1}(T)^{-1}=\sigma_{1}(u) I_{2},
\end{gather*}
so
\begin{gather*}
\phi(S)u=\sigma_{1}(u) S, \qquad \text{with} \quad S:=\sigma_{1}(T) T^{-1} \in \operatorname{GL}_{2}(\operatorname{M}_{\Lambda''}).
\end{gather*}
It follows that there exists $g_{\sigma_{1}}\in \operatorname{M}_{\Lambda''}^{\times}$ (namely, one of the non zero
entries of~$S$) such that
\begin{gather*}
\sigma_{1}(u)=\frac{\phi (g_{\sigma_{1}})}{g_{\sigma_{1}}}u.
\end{gather*}
Consider the norm
\begin{gather*}
\operatorname{N}:=\operatorname{N}_{\operatorname{M}_{\Lambda''}|\operatorname{M}_{\Lambda'}}(g_{\sigma_1})=
\prod\limits_{\sigma \in \operatorname{Gal}(\operatorname{M}_{\Lambda''}|\operatorname{M}_{\Lambda'})} \sigma
(g_{\sigma_1}) \in \operatorname{M}_{\Lambda'}^{\times}.
\end{gather*}
We have
\begin{gather*}
\phi (\operatorname{N}) = \prod\limits_{\sigma \in
\operatorname{Gal}(\operatorname{M}_{\Lambda''}|\operatorname{M}_{\Lambda'})} \sigma \left(\frac{\sigma_{1}(u)
g_{\sigma_{1}}}{u} \right) = \prod\limits_{\sigma \in
\operatorname{Gal}(\operatorname{M}_{\Lambda''}|\operatorname{M}_{\Lambda'})} \sigma(g_{\sigma_{1}}) = \operatorname{N},
\end{gather*}
so $\operatorname{N}=c \in (K^{\phi})^{\times} = \mathbb C^{\times}$.
Up to replacing $g_{\sigma_1}$ by $g_{\sigma_1}c^{-1/k}$, we may assume that
\begin{gather*}
\operatorname{N}_{\operatorname{M}_{\Lambda''}|\operatorname{M}_{\Lambda'}}(g_{\sigma_1})=1.
\end{gather*}
Hilbert's~90 theorem~\cite[Section~X.1]{Ser} ensures that there exists $m \in \operatorname{M}_{\Lambda''}^{\times}$ such
that
\begin{gather*}
{g_{\sigma_1}=\frac{m}{\sigma_1(m)}}.
\end{gather*}
For any $\sigma=\sigma_{1}^{j} \in \operatorname{Gal}(\operatorname{M}_{\Lambda''}|\operatorname{M}_{\Lambda'})$, we set
\begin{gather*}
g_\sigma:= g_{\sigma_1} \sigma_1 (g_{\sigma_1}) \cdots \sigma_1^{j-1} (g_{\sigma_1})=m/\sigma (m)\in
\operatorname{M}_{\Lambda''}^{\times};
\end{gather*}
we have
\begin{gather*}
\sigma(u)=\frac{\phi (g_{\sigma})}{g_{\sigma}}u.
\end{gather*}
It follows that
\begin{gather*}
\widetilde u:= \frac{\phi (m)}{m} u
\end{gather*}
is invariant under the action of $\operatorname{Gal}(\operatorname{M}_{\Lambda''}|\operatorname{M}_{\Lambda'})$ and
hence belongs to $\operatorname{M}_{\Lambda'}^{\times}$.
We have
\begin{gather*}
\phi\left(T'\right)A\left(T'\right)^{-1}= \widetilde{u} I_{2}, \qquad \text{with} \quad T':=mT \in
\operatorname{GL}_{2}(\operatorname{M}_{\Lambda''}).
\end{gather*}
Applying $\sigma \in \operatorname{Gal}(\operatorname{M}_{\Lambda''}|\operatorname{M}_{\Lambda'})$ to this equality, we
get
\begin{gather*}
\phi\left(\sigma(T')\right)A\left(\sigma(T')\right)^{-1}= \widetilde{u} I_{2}.
\end{gather*}
It follows that the matrix
\begin{gather*}
C_{\sigma}:=T' \sigma (T' )^{-1} \in \operatorname{GL}_{2}(\operatorname{M}_{\Lambda''})
\end{gather*}
satisf\/ies $\phi (C_{\sigma}) = C_{\sigma}$ and, hence, that its entries belong to~$K^{\phi} = \mathbb C$.
Moreover, $\sigma \mapsto C_{\sigma}$ is a~$1$-cocyle for the natural action of
$\operatorname{Gal}(\operatorname{M}_{\Lambda''}|\operatorname{M}_{\Lambda'})$ on $\operatorname{GL}_{2}(\mathbb C)$
but this action is trivial so $\sigma \mapsto C_{\sigma}$ is a~group morphism
from~$\operatorname{Gal}(\operatorname{M}_{\Lambda''}|\operatorname{M}_{\Lambda'})$ to $\operatorname{GL}_{2}(\mathbb
C)$.
Since $\operatorname{Gal}(\operatorname{M}_{\Lambda''}|\operatorname{M}_{\Lambda'})$ is cyclic, this implies that there
exists $P\in \operatorname{GL}_{2}(\mathbb C)$ such that, for all $\sigma\in
\operatorname{Gal}(\operatorname{M}_{\Lambda''}|\operatorname{M}_{\Lambda'})$, the matrix
${D_{\sigma}:=P^{-1}C_{\sigma}^{-1}P} \in \operatorname{GL}_{2}(\mathbb C)$ is diagonal.
We have, for all $\sigma \in\operatorname{Gal}(\operatorname{M}_{\Lambda''}|\operatorname{M}_{\Lambda'})$,
\begin{gather*}
\sigma(T'')=D_{\sigma}T'',\qquad  \text{where} \quad T''=(t''_{i,j})_{1 \leq i,j \leq 2}:=P^{-1}T' \in
\operatorname{GL}_{2}(\operatorname{M}_{k \Lambda}).
\end{gather*}
It follows that $u_{1}:=\frac{-t''_{11}}{t''_{12}}$ and $v_{1}:=\frac{-t''_{21}}{t''_{22}}$ are invariant by the action
of $\operatorname{Gal}(\operatorname{M}_{\Lambda''}|\operatorname{M}_{\Lambda'})$ and hence belong to
$\operatorname{M}_{\Lambda'}$.
But $u_{1}$ and $v_{1}$ are solutions of the Riccati equation~\eqref{eq4} (this was already used in the proof of
assertion (2) of Lemma~\ref{lem1}).
So $u_{1}$ and $v_{1}$ are solutions in $\operatorname{M}_{\Lambda'}$ of the Riccati equation~\eqref{eq4}.
\end{proof}

We now come to the main result of this subsection.

\begin{Theorem}
\label{theo2}
The following statements hold:
\begin{enumerate}
\itemsep=0pt
\item[$1.$] The Galois group~$G$ is reducible if and only if the Riccati equation~\eqref{eq4} has at least one solution
in~$\operatorname{M}_{2\Lambda}$.
\item[$2.$] The Galois group~$G$ is completely reducible if and only if the Riccati equation~\eqref{eq4} has at least two
solutions in~$\operatorname{M}_{2\Lambda}$.
\end{enumerate}
\end{Theorem}

\begin{proof}
In virtue of Lemma~\ref{lem1}, it is suf\/f\/icient to prove that:
\begin{enumerate}\itemsep=0pt
\item[(a)] If the Riccati equation~\eqref{eq4} has a~unique solution in~$K$, then it belongs to
$\operatorname{M}_{\Lambda}$.

\item[(b)] If the Riccati equation~\eqref{eq4} has exactly two solutions in~$K$, then they belong to
$\operatorname{M}_{2\Lambda}$.

\item[(c)] If the Riccati equation~\eqref{eq4} has at least three solutions in~$K$, then the Riccati equation~\eqref{eq4}
has at least two solutions in $\operatorname{M}_{2\Lambda}$.
\end{enumerate}

 (a) Assume that the Riccati equation~\eqref{eq4} has a~unique solution~$u$ in~$K$.
Since $u(z)$, $u(z+1)$ and $u(z+\tau)$ are solutions of~\eqref{eq4}, we get
\begin{gather*}
u(z)=u(z+1)=u(z+\tau)
\end{gather*}
and hence $u\in\operatorname{M}_{\Lambda}$.

(b) Assume that the Riccati equation~\eqref{eq4} has exactly two solutions in~$K$ and let $u\in K$ be one of these
solutions.
Since $u(z)$, $u(z+1)$ and $u(z+2)$ are solutions of~\eqref{eq4}, we get $ u(z+2)=u(z).
$ Similarly, we have $u(z+2\tau)=u(z)$.
So $u\in\operatorname{M}_{2\Lambda}$.

(c) What follows is inspired by~\cite[Theorem~4.2]{Hen98}, but is a~little bit subtler.
Assume that the Riccati equation~\eqref{eq4} has at least three solutions in~$K$.
According to Lemma~\ref{lem ric irred bis}, there exist $u\in K$ and ${T=(t_{i,j})_{1 \leq i,j \leq 2} \in
\operatorname{GL}_{2}(K)}$ such that
\begin{gather}
\label{et une equation une}
\phi(T)AT^{-1}=u I_{2}.
\end{gather}
Let $k \in \mathbb N^{*}$ be such that the entries of~$T$ and~$u$ belong to $\operatorname{M}_{k \Lambda}$.
Consider the following f\/ield extensions:
\begin{gather*}
\operatorname{M}_{\Lambda} \subset L \subset \operatorname{M}_{k\Lambda}, \qquad \text{with}\quad L:=\operatorname{M}_{\mathbb
Z+k\tau\mathbb Z}.
\end{gather*}
Applying Lemma~\ref{lem descente} to the extension $\operatorname{M}_{k\Lambda}\vert L$ and to the equation~\eqref{et
une equation une}, we get that the Riccati equation~\eqref{eq4} has two distinct solution $u_{1}$ and $v_{1}$ in~$L$.
If both of them belong to $\operatorname{M}_{2\Lambda}$ then the proof is completed.
Otherwise, up to renumbering, we can assume that $u_{1} \not \in \operatorname{M}_{2\Lambda}$, i.e., that $u_{1}$ is
not $2\tau$-periodic.
Then
\begin{gather*}
u_{1}(z),u_{2}(z):=u_{1}(z+\tau) \qquad \text{and} \qquad u_{3}(z):=u_{1}(z+2\tau)
\end{gather*}
are distinct solutions in~$L$ of the Riccati equation.
For all integers $i,j \in \{1,2,3\}$ with $i<j$ we set $T_{i,j}:=\frac{1}{u_{i}-u_{j}}
\begin{pmatrix}
-u_{j}&1
\\
-u_{i}&1
\end{pmatrix}
\in \operatorname{GL}_{2}(L)$ and we have
\begin{gather*}
\phi (T_{i,j} )A (T_{i,j} )^{-1}=
\begin{pmatrix}
u_{i}&0
\\
0&u_{j}
\end{pmatrix}
\end{gather*}
(this was already used in the proof of assertion~(2) of Lemma~\ref{lem1}).
Therefore,
\begin{gather*}
\phi\big(T_{1,3} (T_{1,2} )^{-1}\big)
\begin{pmatrix}
u_{1}&0
\\
0&u_{2}
\end{pmatrix}
=
\begin{pmatrix}
u_{1}&0
\\
0&u_{3}
\end{pmatrix}
T_{1,3} (T_{1,2} )^{-1}.
\end{gather*}
Equating the second columns in this equality, we see that there exists $f \in L^{\times}$ such that either
$u_{1}=\frac{\phi f}{f} u_{2}$ or $u_{3}=\frac{\phi f}{f} u_{2}$; up to renumbering, we may assume that the former
equality holds true.
Then, we have
\begin{gather}
\label{et une deux equa}
\phi\big(\widetilde{T}\big)A\widetilde{T}^{-1}= u_{1} I_{2}
\end{gather}
with
\begin{gather*}
u_{1} \in L^{\times} \qquad \text{and} \qquad \widetilde{T}:=
\begin{pmatrix}
1&0
\\
0&f
\end{pmatrix}
T_{1,2}\in \operatorname{GL}_{2}(L).
\end{gather*}
Applying Lemma~\ref{lem descente} to the extension $L \vert \operatorname{M}_{\Lambda}$ and to the equation~\eqref{et
une deux equa}, we see that the Riccati equation~\eqref{eq4} has $2$ distinct solutions in $\operatorname{M}_{\Lambda}$.
This concludes the proof.
\end{proof}
